\scanpage{8}\mspage{7}
Je voudrais maintenant faire passer les axiomes de disjonction-compatibilité avant ceux sur les lieux, pour ne pas \inserted{\tikz[baseline=(t.base),x=.6em,y=.6em]{\node[inner sep=0pt,anchor=west] (t) at (.15,.55) {avoir l'air de}; \draw (0,-.3)..controls(.1,.05)and(.6,.15)..(6.8,.25)..controls(7.0,.3)and(7.1,.45)..(7.1,.6);}} les en faire dépendre. Les axiomes sont ``facultatifs'' pour un atelier, et il y aura pas mal \unsure{des'|des|d'} implications mutuelles avec les axiomes sur les lieux.

Une notion essentielle ici est celle de ``$F$ et $G$ \uline{disjoints}''
\[
 F\parallel G\ \overset{\mathrm{dfn}}{\Longleftrightarrow}\ F\mathrel{\tikz[baseline=-.4ex,x=.65em,y=.65em]{\draw (0,1.1)--(1,.5)--(0,-.1) (1,1.1)--(0,.5)--(1,-.1);}}G\quad\text{et}\quad F\cap G=\varnothing_F
\]
notion qui ne dépend que de $\leq$. On a d'ailleurs
\[
 F\parallel G\ \Longleftrightarrow\ \forall X\in\widetilde F,\ Y\in\widetilde G,\quad X\parallel Y
\]
donc interprétation en \unsure{termes|terme} de $X\parallel Y$ dans $M$, qui y signifie
\[
\begin{aligned}
X\mathrel{\underset{M}{\parallel}}Y\ \Longleftrightarrow\
 X\mathrel{\underset{M}{\mathrel{\tikz[baseline=-.4ex,x=.65em,y=.65em]{\draw (0,1.1)--(1,.5)--(0,-.1) (1,1.1)--(0,.5)--(1,-.1);}}}}Y\quad &\text{et}\quad\{X,Y\}\ \text{n'est \uline{pas}}\\
&\text{minoré dans } M,\leq\\
&\text{i.e. }\nexists Z\ \text{avec }Z\leq X,Z\leq Y
\end{aligned}
\]

% The shared prefix résul- is clear. Only the t/ending cluster is overwritten; the former terminal group cannot be independently recovered. The surviving ending favors te.
% The curved stroke beneath précédents descends into the axiomes / que gap: an insertion guide, not an underline.
(\uline{NB} Il résul\deleted{\illegible{1}}\unsure{te|tera} des axiomes \inserted{\tikz[baseline=(t.base),x=.6em,y=.6em]{\node[inner sep=0pt,anchor=west] (t) at (.1,.55) {précédents}; \draw (1.2,-.4)..controls(1.3,.05)and(2.4,.15)..(5.6,.35)..controls(6.1,.38)and(6.3,.48)..(6.4,.68);}} que ça implique que $\nexists Z$ avec $Z\ll X,Z\ll Y$, \deleted{i.e. c'est un $\unsure{\exists}$} et \deleted{\unsure{la}} réciproque résultera \unsure{d'axiome|d'axiomes} ci-dessous\ldots)

% The marginal note is slanted. A long diagonal of the general cancellation intersects its lower/right portion; no separate note-only cancellation is established.
% Three long diagonals cancel the following At DC 1 and Cor draft as a block; the outer deleted macros record that state, distinct from the earlier local cancellations.
% slanted
\margin{Un atelier sera dit \uline{\unsure{spacieux|specieux}} s'il satisfait cette condition. Contre-ex. : les ateliers simpliciaux\ldots}%
\deleted{\uline{At DC} 1 \quad (\uline{Axiome des ateliers \unsure{spacieux|specieux}}) Soient $X,Y\in M$, \inserted{soit} \deleted{\unsure{et}} $K=X\cap Y=\underset{Z\in\widetilde X\cap\widetilde Y}{\mathrm{Sup}}\,Z$. \deleted{\unsure{Il}} Supposons que \inserted{pour $X',Y'\in M$}}
\[
\deleted{\deleted{\forall}\ X'\ll X,\ X'\parallel K,\ Y'\ll Y,\ Y'\parallel K\ \Longrightarrow\ X'\parallel Y'.}
\]
\deleted{Alors [NB et alors seulement\ldots\ \inserted{cf Cor}] $X\mathrel{\tikz[baseline=-.4ex,x=.65em,y=.65em]{\draw (0,1.1)--(1,.5)--(0,-.1) (1,1.1)--(0,.5)--(1,-.1);}}Y$.}

\deleted{\uline{Cor} Soient $F,G\in\mathscr F$, $K=F\cap G=\underset{Z\in\deleted{\unsure{\widetilde X|X}}\unsure{F|\widetilde F}\cap\deleted{\unsure{\widetilde Y|Y}}\widetilde G}{\mathrm{Sup}}\,Z$.}
\[
\begin{aligned}
&\deleted{\text{Alors } F\mathrel{\tikz[baseline=-.4ex,x=.65em,y=.65em]{\draw (0,1.1)--(1,.5)--(0,-.1) (1,1.1)--(0,.5)--(1,-.1);}}G\quad \text{\tikz[baseline=(s.base)]{\node[inner sep=1pt] (s) {ses};\draw (s.center) ellipse [x radius=1.05em,y radius=.72em];}}\quad \forall X,\deleted{\unsure{\varphi|y}}Y\in M}\\
&\deleted{X\ll F,\ Y\ll G,\ X\parallel K,\ Y\parallel K\ \Longrightarrow\ X\parallel Y}
\end{aligned}
\]
